% BNSF Data Scientist I/II (95561): adapted from the existing 3-page comprehensive industry resume; retain layout and verified experience.
% Python replaces template tokens. Edit layout here, content in JSON/portfolio JS.
\documentclass[9pt,letterpaper]{extarticle}
\usepackage[left=0.68in,right=0.68in,top=0.45in,bottom=0.45in]{geometry}
\usepackage[T1]{fontenc}
\input{glyphtounicode}
\pdfgentounicode=1
\usepackage[utf8]{inputenc}
\usepackage{lmodern}
\usepackage{microtype}
\usepackage{pifont}
\usepackage{xcolor}
\usepackage{enumitem}
\usepackage{titlesec}
\usepackage[hidelinks]{hyperref}
\usepackage{tabularx,array,parskip,ragged2e,needspace}
\pagestyle{empty}
\setlength{\parindent}{0pt}
\setlength{\parskip}{1pt}
\definecolor{accent}{HTML}{155B8A}
\titleformat{\section}{\large\bfseries\color{accent}}{}{0pt}{}[\vspace{0.5pt}\color{accent}\titlerule]
\titlespacing*{\section}{0pt}{3pt}{2pt}
% Compact subsections beneath the main accent-rule sections.
\titleformat{\subsection}{\normalsize\bfseries\color{accent}}{}{0pt}{}
\titlespacing*{\subsection}{0pt}{3pt}{1pt}
\setlist[itemize]{leftmargin=1.15em,itemsep=0.8pt,topsep=1pt,parsep=0pt,partopsep=0pt}
\newcommand{\role}[4]{%
  \Needspace{4\baselineskip}\textbf{#1}\hfill\textbf{#2}\\[-1pt]
  \textit{#3}\hfill\textit{#4}\par}
\newcommand{\entry}[2]{\Needspace{4\baselineskip}\textbf{#1}\hfill\textit{#2}\par}
\newcommand{\skillrow}[2]{\textbf{#1:} & #2\\}
\hypersetup{pdftitle={Waseem Raza, PhD - BNSF Data Scientist I II - Job 95561},pdfauthor={Waseem Raza, PhD}}
\begin{document}

\begin{center}
{\LARGE\bfseries Waseem Raza, PhD}{\normalsize\mdseries }\\[2pt]
{\large\bfseries Data Scientist \textbar{} Machine Learning, Predictive Analytics \& Cloud Data}\\[1pt]
{\small Python \textbar{} SQL \textbar{} Azure Databricks \textbar{} Statistical Modeling \textbar{} Operational Analytics}\\[3pt]
{\footnotesize \ding{41}\enspace \href{mailto:waseem.placements@gmail.com}{waseem.placements@gmail.com} \textbar{} \href{mailto:waseemraza844@gmail.com}{waseemraza844@gmail.com}\\[1pt]New Jersey, United States \quad \ding{37}\enspace \href{tel:+17328039703}{+1 732-803-9703}\\[1pt]\href{https://www.linkedin.com/in/waseem-raza-phd-9a737858/}{LinkedIn} \textbar{} \href{https://github.com/WaseemRaza844}{GitHub} \textbar{} \href{https://scholar.google.com/citations?user=d58Drr0AAAAJ\&hl=en}{Google Scholar} \textbar{} \href{https://waseemraza844.github.io/portfolio/learning.html}{Coursera Courses} \textbar{} \href{https://waseemraza844.github.io/portfolio/}{Portfolio}}
\end{center}
\vspace{-5pt}
\Needspace{3\baselineskip}
\section{Professional Profile}
Data Scientist and Machine Learning Engineer with a PhD in Electrical and Computer Engineering and 7+ years of combined applied research, engineering, and data-driven problem-solving experience. Deliver statistical and machine-learning analysis of complex operational data, from extracting and validating records to feature engineering, model evaluation, and stakeholder communication. At AT\&T, improved automated E911 root-cause classification performance from approximately 20\% to 80\%; at Impac, built Python/Azure data workflows that reduced data-preparation time by 40\%. Hands-on with Python, SQL, scikit-learn, PyTorch, Azure Databricks, Synapse, Azure Data Factory, model validation, anomaly detection and cross-domain analytics. Interested in applying these methods to freight operations, asset reliability, and process improvement.
\Needspace{2\baselineskip}
\subsection{Technical Competencies}
\Needspace{2\baselineskip}\noindent{\small\textbf{Data Science \& Statistics:} Python, SQL, Pandas, NumPy, SciPy, feature engineering, exploratory analysis, statistical testing, regression, classification, clustering, anomaly detection, time-series analysis, cross-validation, error analysis and experimental design}\par
\Needspace{2\baselineskip}\noindent{\small\textbf{ML \& AI Frameworks:} scikit-learn, PyTorch, TensorFlow/Keras, LightGBM, XGBoost, neural networks, transfer learning, model selection, imbalanced learning, interpretability, model robustness and domain adaptation}\par
\Needspace{2\baselineskip}\noindent{\small\textbf{Cloud Data Engineering:} Azure Databricks, Apache Spark/PySpark, Azure Data Factory, Synapse Analytics, Azure Functions, Python/SQL ETL and ELT, ingestion, transformations, data quality, validation, pipeline troubleshooting and Azure DevOps CI/CD}\par
\Needspace{2\baselineskip}\noindent{\small\textbf{Visualization \& Reporting:} Matplotlib, Plotly, Pandas, Jupyter, analytical reports, operational diagnostics, metrics and communicating model findings to technical and nontechnical stakeholders}\par
\Needspace{2\baselineskip}\noindent{\small\textbf{Software \& MLOps:} Git, GitHub Actions, Linux, Docker, REST APIs, reproducible experiments, model evaluation, deployment-oriented workflows, basic Grafana exposure and systematic troubleshooting}\par
\Needspace{2\baselineskip}\noindent{\small\textbf{Applied Analytics \& Optimization:} Network event data, reliability analysis, anomaly detection, large-scale logs, simulation, optimization, data-efficient learning and decision support for operational systems}\par
\Needspace{2\baselineskip}\noindent{\small\textbf{Generative AI:} LLM-enabled analytics, retrieval-augmented generation (RAG), embeddings, LangChain, LangGraph and structured extraction; applied research and portfolio workflows}\par
\Needspace{4\baselineskip}\textbf{Delivery and Collaboration}\par
\begin{itemize}
\item \textbf{Business problem framing:} Translate ambiguous operational questions into measurable analytical objectives, defensible baselines and actionable results.
\item \textbf{Cross-functional collaboration:} Work with network and software engineers, geologists, researchers and service-assurance stakeholders to define data requirements and explain findings.
\item \textbf{Reliable experimentation:} Prioritize reproducible data processing, data-quality checks, cross-validation, failure analysis and clear interpretation of model limitations.
\item \textbf{Communication:} Summarize complex modeling choices and results for mixed technical audiences through reports, plots and documented comparisons.
\end{itemize}
\Needspace{3\baselineskip}
\section{Professional Experience}
\role{Machine Learning Engineer}{02/2025 - Present}{AT\&T, contractor via Innovatix Inc.}{Bedminster, NJ}
\begin{itemize}\interlinepenalty=10000
\item Develop and evaluate ML models for critical E911 service-assurance event data, improving automated root-cause classification performance from approximately 20\% to 80\% and enabling faster analysis of operational issues.
\item Engineered text-feature and domain-adaptive classification workflows, handling label imbalance, data-quality differences and distribution shifts across operational domains.
\item Analyzed work-order descriptions, network context and location-linked records to identify failure patterns; explored GenAI-assisted diagnostics in service-assurance prototypes.
\item Built repeatable model-validation and error-analysis procedures to compare precision, recall, failure modes and cross-domain generalization for high-impact workflows.
\item Collaborated with network engineers, software teams and public-safety stakeholders to turn operational questions into measurable ML tasks and communicate data-backed findings.
\end{itemize}
\role{Data Scientist and AI/ML Engineer}{09/2024 - 12/2024}{Impac Exploration Services}{Weatherford, OK}
\begin{itemize}\interlinepenalty=10000
\item Built an LLM-enabled analytical pipeline for extracting and interpreting information from geochemical well data.
\item Developed Python/SQL-supported ingestion, cleansing, transformation and validation workflows with Azure Databricks, Azure Functions and Synapse Analytics.
\item Automated ETL processing for large exploration datasets, reducing data-preparation time by 40\% and improving reproducibility.
\item Implemented cross-validation, data-quality, and consistency checks that improved reliability of analytical outputs by 20\%.
\item Collaborated with geologists to convert domain requirements and technical information into AI-ready datasets and actionable analytical outputs.
\end{itemize}
\role{Advanced-Degrees Intern}{06/2024 - 08/2024}{AT\&T Labs}{Austin, TX}
\begin{itemize}\interlinepenalty=10000
\item Analyzed more than 500K LTE/5G diagnostic logs with statistical and ML methods, identifying anomalous patterns and reducing diagnostic time by 20\%.
\item Investigated 5G NR-NTN cell-selection behavior and improved decision algorithms by 15\% through data-driven analysis and validation.
\item Conducted structured analysis of large operational datasets and communicated technical findings to engineering teams supporting next-generation wireless systems.
\end{itemize}
\role{Advanced Wireless Technology Development Intern}{09/2022 - 01/2023}{Futurewei Technologies}{Schaumburg, IL}
\begin{itemize}\interlinepenalty=10000
\item Improved autoencoder-based channel-state-information compression by 25\% using transfer learning and systematic experimentation on large wireless datasets.
\item Conducted 50+ experiments to validate model performance, robustness, and generalization under changing data conditions.
\item Collaborated with research and engineering stakeholders to translate modeling results into recommendations for next-generation communication systems.
\end{itemize}
\role{Advanced-Degrees Intern}{06/2022 - 08/2022}{AT\&T Labs}{Bedminster, NJ}
\begin{itemize}\interlinepenalty=10000
\item Designed ML and statistical methods for analyzing users' timing in 5G standalone networks, improving analytical precision by 20\%.
\item Analyzed large-scale user trace data with SQL and visualization tools and communicated findings to senior technical leadership.
\item Converted raw operational traces into interpretable performance indicators to support diagnosis and network optimization.
\end{itemize}
\role{Graduate Research Assistant}{09/2019 - 07/2025}{AI4Networks, University of Oklahoma}{Tulsa, OK}
\begin{itemize}\interlinepenalty=10000
\item Designed reusable Python simulation and data-generation platforms for large-scale experimentation, including a 3GPP-compliant network simulator and an AI-enabled 5G testbed.
\item Built transfer-learning, adversarial-learning, GAN, autoencoder, and reinforcement-learning pipelines for sparse-data and cross-domain modeling.
\item Developed robustness evaluation frameworks for models operating under distribution shifts, missing data, and limited labeled samples; led 100+ experiments from problem formulation through validation.
\item Developed scalable self-healing and outage-management frameworks that improved modeled network resilience by 30\%.
\item Published peer-reviewed research, presented technical results internationally, and mentored junior researchers in ML implementation, experimental design, and technical writing.
\end{itemize}
\Needspace{3\baselineskip}
\section{Education}
\entry{PhD, Electrical and Computer Engineering}{09/2019 - 07/2025}
University of Oklahoma, Norman, Oklahoma, USA\par
\entry{MSc, Telecommunication Engineering}{08/2014 - 08/2016}
University of Engineering and Technology, Taxila, Pakistan\par
\entry{BSc, Telecommunication Engineering}{11/2010 - 06/2014}
University of Engineering and Technology, Taxila, Pakistan\par
\Needspace{3\baselineskip}
\section{Selected Projects \& Research}
\Needspace{2\baselineskip}
\subsection{Projects \& Applied Research}
\entry{AI-Ready Lakehouse Pipeline}{Portfolio Project \textbar{} 2024}
Designed an AI-ready lakehouse pipeline that prepares structured and semi-structured data for analytics and downstream machine-learning workloads. Implemented repeatable ingestion, cleansing, transformation, validation, and scalable processing patterns suitable for reproducible cloud data workflows.\par
{\small\textit{Tools \& Methods:} Azure Databricks, Apache Spark, PySpark, ETL/ELT, Data Ingestion, Schema Transformation, Data Validation, Feature Preparation, Cloud Data Engineering\par}
\vspace{1pt}
\entry{Financial ML and Forecasting}{Portfolio Project \textbar{} 2024-2025}
Developed machine-learning workflows for structured credit-risk and loan-default datasets, together with time-series forecasting experiments. Applied data preparation, feature engineering, model comparison, cross-validation, and predictive-performance analysis to structured financial datasets.\par
{\small\textit{Tools \& Methods:} Python, scikit-learn, Classification, Credit Risk, Loan Default Modeling, Time-Series Forecasting, Feature Engineering, Cross-Validation, Model Evaluation\par}
\vspace{1pt}
\entry{Wearable Intelligence for Health Monitoring}{Research Project \textbar{} 2024-2026}
Developed deep-learning workflows for recognizing medication-taking and dietary activities from smartwatch accelerometer and gyroscope signals in free-living settings. Evaluated windowing, feature extraction, CNN/LSTM-style sequence models, cross-user generalization, and recall-oriented performance for health-monitoring applications.\par
{\small\textit{Tools \& Methods:} Wearable AI, Smartwatch IMU, Time Series, Deep Learning, CNN/LSTM, Feature Engineering, Cross-User Evaluation, Activity Recognition, Healthcare AI\par}
\vspace{1pt}
\entry{\href{https://github.com/WaseemRaza844/agentic-ai-projects-by-was}{Agentic AI and RAG Applications}}{Portfolio Project \textbar{} 2026}
Built grounded GenAI workflows that combine document retrieval, vector search, stateful orchestration, tool-using agents, memory, evaluation, and multi-agent handoffs. Emphasized controllable execution, traceable context flow, and reliable responses for multi-step tasks rather than single unconstrained model calls.\par
{\small\textit{Tools \& Methods:} Python, LangGraph, LangChain, CrewAI, AutoGen, RAG, Vector Databases, Embeddings, Tool Calling, Agent Memory, Evaluation\par}
\vspace{1pt}
\entry{Propagation Models \& Digital Twins}{Research Project \textbar{} 2024}
Developed robust data-driven radio-propagation models for wireless digital twins using domain-informed generative learning and scarce measured data. Combined 3GPP-informed pretraining, GAN-based augmentation, domain adaptation, and statistical validation to improve resilience under distribution shift.\par
{\small\textit{Tools \& Methods:} Propagation Modeling, Digital Twins, GANs, Domain Adaptation, 3GPP Models, Synthetic Data, Statistical Validation, Wireless AI, Distribution Shift\par}
\vspace{1pt}
\entry{SyntheticNET: 3GPP-Compliant Network Simulation}{Research Project \textbar{} 2020-2024}
Developed a Python-based, 3GPP-aligned simulation environment for realistic mobility, propagation, and cellular-network experimentation when measured data is sparse or costly. Supported controlled data generation and repeatable studies in propagation modeling, network optimization, digital twins, and resilient wireless AI.\par
{\small\textit{Tools \& Methods:} Python, 3GPP, Wireless Simulation, Mobility Modeling, Propagation Modeling, Synthetic Data, Experiment Design, Digital Twins, Network Analytics\par}
\vspace{1pt}
\entry{\href{https://github.com/WaseemRaza844/system-design-masterclass}{System Design Case-Study Portfolio}}{Guided Project \textbar{} 2026}
Developed architecture case studies that translate ambiguous product requirements and traffic assumptions into service boundaries, APIs, data models, storage choices, caching strategies, and reliability targets. Documented scalability, consistency, capacity, and failure-mode tradeoffs to support production-oriented system design decisions.\par
{\small\textit{Tools \& Methods:} System Design, REST APIs, Distributed Systems, Data Modeling, Caching, Scalability, Reliability, Capacity Planning, Architecture Tradeoffs\par}
\vspace{1pt}
\Needspace{2\baselineskip}
\subsection{Publications}
\entry{Real-World Medication Adherence Monitoring Using Hybrid Neural Networks and Smartwatches}{2026}
\textit{HealthINF \textbar{} Accepted}\par
\entry{\href{https://doi.org/10.3390/s25247651}{On Multi-Parameter Optimization and Proactive Reliability in 5G and Beyond Cellular Networks}}{2025}
\textit{Sensors 2025}\par
\entry{\href{https://doi.org/10.1109/PIMRC59610.2024.10817403}{An AI-Driven Framework for Enhancing Resilience in Propagation Models to Enable Digital Twin}}{2024}
\textit{IEEE PIMRC}\par
\entry{Data-Driven Intelligent Outage Management for High Shadowing Environments in 5G\&B Networks}{2024}
\textit{IEEE SmartNets}\par
\entry{\href{https://arxiv.org/abs/2203.15899}{Toward a Hybrid RF/Optical Lunar Communication System (LunarComm)}}{2022}
\textit{IEEE Network}\par
\entry{Towards Positioning Error Impact Characterization and Minimization in User-Centric RAN}{2022}
\textit{IEEE WCNC}\par
\entry{\href{https://doi.org/10.1016/j.aeue.2020.153244}{Unification of RF Energy Harvesting Schemes under Mixed Rayleigh-Rician Fading Channels}}{2020}
\textit{AEÜ - International Journal of Electronics and Communications}\par
\Needspace{3\baselineskip}
\section{Certifications \& Professional Development}
\Needspace{3\baselineskip}\noindent\textbf{1. Generative AI Applications with Amazon Bedrock} {\small --- \textcolor{blue}{\underline{\href{https://www.coursera.org/account/accomplishments/verify/3W16P8EE4BJZ}{Coursera}}} \textbar{} Completed \textbar{} August 27, 2026 \textbar{} \textcolor{blue}{\underline{\href{https://waseemraza844.github.io/portfolio/certificates/genai/aws-bedrock/c2-generative-ai-applications-bedrock.pdf}{Certificate}}}}\par
\textbf{Focus:} Covers development of generative-AI applications on Amazon Bedrock, including foundation-model access, prompt-based application logic, model interaction, and managed AWS services. Emphasizes practical cloud patterns for building and operating LLM-powered applications.\par
{\small\textbf{Tools \& Skills:} Amazon Bedrock, foundation models, prompt engineering, AWS generative-AI services, LLM application development, cloud AI workflows\par}
\vspace{2pt}
\Needspace{3\baselineskip}\noindent\textbf{2. Agentic AI with LangChain and LangGraph} {\small --- \textcolor{blue}{\underline{\href{https://www.coursera.org/account/accomplishments/verify/L4EY4SLGCPUE}{Coursera}}} \textbar{} Completed \textbar{} February 9, 2026}\par
\textbf{Focus:} Covers design of stateful agentic applications using LangChain and LangGraph, including graph-based orchestration, tool use, memory and state, routing, and multi-step workflows. Emphasizes controllable agent execution and composition of reliable LLM-driven processes.\par
{\small\textbf{Tools \& Skills:} LangChain, LangGraph, tool-calling, state and memory, workflow orchestration, routing, agent design\par}
\vspace{2pt}
\Needspace{3\baselineskip}\noindent\textbf{3. Machine Learning} {\small --- \textcolor{blue}{\underline{\href{https://www.coursera.org/account/accomplishments/verify/VU3Y5TBLVLZT}{Stanford Online / Coursera}}} \textbar{} Completed \textbar{} March 29, 2020}\par
\textbf{Focus:} Covers foundational machine-learning methods for supervised and unsupervised problems, including regression, classification, model evaluation, regularization, clustering, and practical model-selection considerations. Emphasizes translating data into predictive models while controlling generalization error.\par
{\small\textbf{Tools \& Skills:} supervised learning, regression, classification, regularization, clustering, model evaluation, bias-variance tradeoffs, machine-learning workflows\par}
\vspace{2pt}
\Needspace{3\baselineskip}\noindent\textbf{4. Introduction to Data Analytics in Google Cloud} {\small --- \textcolor{blue}{\underline{\href{https://www.coursera.org/account/accomplishments/verify/R9FGPRZ3NJE8}{Coursera}}} \textbar{} Completed \textbar{} July 24, 2024}\par
\textbf{Focus:} Introduces cloud-native data analytics in Google Cloud, covering the analytics lifecycle, data sources, storage and processing concepts, and how cloud services support scalable analysis. Emphasizes organizing data work to produce reliable analytical insight.\par
{\small\textbf{Tools \& Skills:} Google Cloud, cloud analytics, data lifecycle, data-ingestion concepts, storage and processing, analytical workflows, data-driven decision support\par}
\vspace{2pt}
\Needspace{3\baselineskip}\noindent\textbf{5. Introduction to Deep Learning \& Neural Networks with Keras} {\small --- \textcolor{blue}{\underline{\href{https://www.coursera.org/account/accomplishments/verify/Y6NP4GWLS4DZ}{IBM / Coursera}}} \textbar{} Completed \textbar{} August 12, 2026}\par
\textbf{Focus:} Introduces deep-learning and neural-network fundamentals using Keras, covering network architecture, training, validation, and model improvement. Connects concepts such as activation functions and optimization with practical neural-network implementation and evaluation.\par
{\small\textbf{Tools \& Skills:} Keras, deep learning, neural networks, model training, validation, activation functions, optimization, model evaluation\par}
\vspace{2pt}
\end{document}
